\documentclass[10pt,a4paper]{article}
\usepackage[margin=0.8in]{geometry}
\usepackage{amsmath,amssymb,hyperref}
\hypersetup{colorlinks=true,urlcolor=blue,linkcolor=blue}
\setlength{\parindent}{0pt}\setlength{\parskip}{5pt}
\begin{document}
\begin{center}
{\Large Finite exceptional sequences and infinitely many positive points}\\[4pt]
TLMC 00000000553 \quad lizaixi01 \quad 24 October 2026
\end{center}

\textbf{The precise objects.}
We use complete exceptional sequences in the category $\operatorname{mod}\Lambda$
of finite-dimensional modules, counted up to componentwise module isomorphism.
We use ordinary positive rational points of the associated coefficient-free
cluster $\mathcal A$-variety. These are the standard conventions specified below.
For the genuine rank-two algebra
\[
 \Lambda=\mathbb Q\times\mathbb Q,
\]
the first set is finite and nonempty, whereas the second is infinite.
Thus their cardinalities cannot agree. The source includes rank two and
does not impose connectedness. This counterexample uses the disconnected
Dynkin type $A_1\sqcup A_1$.
It does not concern counts of clusters, tropical points, canonical-basis
labels, or exceptional objects with all possible derived shifts.

Fakieh--Ringel explicitly allow disconnected Dynkin diagrams and define
the rank as the number of simple module isomorphism classes.
Their complete sequences have length equal to that rank; for $i>j$ they
require $\operatorname{Hom}(M_i,M_j)=0=\operatorname{Ext}^1(M_i,M_j)$,
and all counts are up to isomorphism \cite{fr}.
Fock--Goncharov define positive points using a positive atlas and show that
each torus chart gives a bijection with the positive points; here the
semifield is $\mathbb Q_{>0}$ \cite{fg}.

\textbf{The full module category, not a list of chosen names.}
Let
\[
 I_0=0\times\mathbb Q,\quad I_1=\mathbb Q\times0,
 \qquad S_i=\Lambda/I_i.
\]
The formal construction starts with all native \texttt{ModuleCat} objects
with finite generation over $\Lambda$, hence finite dimension over $\mathbb Q$.
The ring is semisimple. Every submodule has a complement; therefore an
indecomposable nonzero module is simple. Conversely each simple is
indecomposable.

Every simple module is isomorphic to a quotient $\Lambda/I$ for a
maximal ideal $I$. The classification of ideals of a product of fields
gives exactly $I_0,I_1$. An isomorphism preserves the actual annihilator,
and $\operatorname{Ann}_\Lambda S_i=I_i$, so $S_0,S_1$ are not isomorphic.
The quotient of the full simple-module type by linear isomorphism is
therefore in bijection with two elements. The rank is computed as its
actual cardinality, giving $\operatorname{rank}\Lambda=2$.

Semisimplicity also gives the lifting property for every surjection,
so every module is projective. The native groups
$\operatorname{Ext}^1_\Lambda(M,N)$ are zero for all $M,N$.
Schur's lemma shows that each endomorphism ring of an indecomposable is
a division ring. It also gives
\[
 \operatorname{Hom}_\Lambda(S_0,S_1)=0=
 \operatorname{Hom}_\Lambda(S_1,S_0).
\]
Thus $(S_0,S_1)$ is an actual complete exceptional sequence.
For an arbitrary complete sequence, replacing each entry by its simple
isomorphism class gives an injective map from sequence isomorphism classes
into the finite set of ordered pairs of simple classes. This proves
finiteness without substituting a prescribed integer for the count.

\newpage
\textbf{The associated exchange data.}
The two simple residue fields are both $\mathbb Q$, by the actual projection
maps and the quotient-kernel isomorphism theorem. All extensions between
the simples vanish. Hence the Ext-quiver has two split vertices, no edges,
and the coefficient-free exchange matrix computed from their Ext dimensions is
\[
 \varepsilon_{ij}=
 \dim_\mathbb Q\operatorname{Ext}^1(S_i,S_j)
 -\dim_\mathbb Q\operatorname{Ext}^1(S_j,S_i)=0.
\]
The multipliers are one and there are no frozen vertices. These are the
associated $A_1\sqcup A_1$ seed data, not an unrelated matrix assigned to the
algebra. Matrix mutation preserves this zero exchange matrix for every
finite mutation word.

\textbf{Tori, actual mutations, and global positive points.}
The lattice is $L=\mathbb Z^2$, represented by finitely supported functions
on the two vertices. The rational torus is the actual group of characters
\[
 H(\mathbb Q)=\operatorname{Hom}(L,\mathbb Q^\times).
\]
Evaluating a character on the two basis vectors is proved to give a
bijection with $(\mathbb Q^\times)^2$. Its positive part is consequently
in bijection with $P=(\mathbb Q_{>0})^2$.

For the cluster $\mathcal A$ coordinates, equation (14) of
Fock--Goncharov is the sum of two exchange monomials divided by the
mutated coordinate. The two monomials are actual finite products.
Because every exchange exponent is zero, each product is one, and
mutation is exactly
\[
 \mu_0(x_0,x_1)=(2/x_0,x_1),\qquad
 \mu_1(x_0,x_1)=(x_0,2/x_1).
\]
Both maps are positive rational involutions. We use $\mathcal A$ throughout;
no formula for the different $\mathcal X$ construction is invoked.

The formal atlas retains a separate chart for every finite labelled
mutation word $w$. Composing its mutations gives a bijection
$b_w:P\to P$ to the base chart. The transition between charts $u,v$ is
$b_v^{-1}b_u$, satisfies the identity and cocycle equations, and on an
adjacent pair is the actual mutation above. Global points are the quotient
of pairs $(w,x)$ by
\[
 (u,x)\sim(v,y)\quad\Longleftrightarrow\quad b_u(x)=b_v(y).
\]
This is the positive-point gluing construction: it does not quotient a
single torus by the mutation group. Each chart is proved bijective onto
the resulting global positive-point set, agreeing with the positive atlas
definition in \cite{fg}.

Finally the distinct points $(n+1,1)$ for $n\in\mathbb N$ inject through
the base chart into that global set. It is infinite. A finite set of
complete exceptional-sequence classes cannot be in bijection with it.
Only the printed equality under the stated standard conventions is refuted;
no conventional finite cluster-tilting correspondence is contradicted.

\begin{thebibliography}{9}
\bibitem{fr} W. Fakieh and C. M. Ringel,
\emph{The number of complete exceptional sequences for a Dynkin algebra},
pp. 1, 4. \url{https://www.math.uni-bielefeld.de/~ringel/opus/e2.pdf}.
\bibitem{fg} V. V. Fock and A. B. Goncharov,
\emph{Cluster ensembles, quantization and the dilogarithm},
arXiv:math/0311245v7 (2009), pp. 4--8, especially Definition 1.3
and equation (14). \url{https://arxiv.org/abs/math/0311245}.
\end{thebibliography}
\end{document}
